\DocumentMetadata{
  pdfversion=2.0,
  pdfstandard=ua-2,
  lang=en-US,
  tagging=on,
  tagging-setup={math/setup=mathml-SE,tabsorder=structure}
}
\documentclass[11pt,letterpaper]{article}
\usepackage[margin=0.68in,includefoot]{geometry}
\usepackage{newcomputermodern}
\usepackage{amsmath,amscd,mathtools}
\providecommand{\square}{\mdlgwhtsquare}
\usepackage[hidelinks]{hyperref}
\hypersetup{
  pdftitle={Analysis First-Year Exam, August 2020, Part 2},
  pdfauthor={University of Florida Department of Mathematics},
  pdfsubject={Graduate mathematics examination},
  pdfdisplaydoctitle=true,
  bookmarksopen=true
}
\setcounter{secnumdepth}{0}
\setlength{\parindent}{0pt}
\setlength{\parskip}{0.28em}
\setlength{\emergencystretch}{3em}
\setlength{\leftmargini}{2.15em}
\setlength{\leftmarginii}{2.65em}
\setlength{\leftmarginiii}{3em}
\renewcommand{\labelenumi}{\textbf{\theenumi.}}
\pagestyle{empty}
\raggedbottom
\allowdisplaybreaks
\newenvironment{examproblems}[1]{%
  \begin{enumerate}%
  \setcounter{enumi}{#1}%
  \setlength{\topsep}{0.35em}%
  \setlength{\partopsep}{0pt}%
  \setlength{\itemsep}{0.48em}%
  \setlength{\parsep}{0.18em}%
}{\end{enumerate}}
\newenvironment{examsubparts}{%
  \begin{enumerate}%
  \setlength{\topsep}{0.18em}%
  \setlength{\partopsep}{0pt}%
  \setlength{\itemsep}{0.16em}%
  \setlength{\parsep}{0.1em}%
}{\end{enumerate}}
\newenvironment{examinstructions}{%
  \par\begingroup\itshape
}{%
  \par\endgroup\vspace{0.18em}
}
\makeatletter
\renewcommand\section{\@startsection{section}{1}{0pt}%
  {-1.25ex plus -.25ex minus -.1ex}%
  {0.55ex plus .1ex}%
  {\normalfont\large\bfseries}}
\renewcommand{\@maketitle}{%
  \newpage\null\vskip -1.2em%
  {\footnotesize\bfseries
    \href{https://www.ufl.edu/}{University of Florida}, \href{https://math.ufl.edu/}{Department of Mathematics}\par}%
  \vskip 0.18em%
  \tagstructbegin{tag=Title}%
    {\LARGE\bfseries\href{https://gma.math.ufl.edu/exams/analysis-first-year/}{Analysis First-Year Exam}, August 2020, Part 2\par}%
  \tagstructend%
  \vskip 0.35em%
  \tagmcbegin{artifact}%
    \hrule height 1pt%
  \tagmcend%
  \vskip 0.55em%
}
\makeatother
\title{Analysis First-Year Exam, August 2020, Part 2}
\author{}
\date{}
\begin{document}
\maketitle
\thispagestyle{empty}
\begin{examinstructions}
Answer FOUR questions in detail. State carefully any results used without proof.
\end{examinstructions}
\begin{examproblems}{0}
\item
The function \(f:[0,1]\to\mathbb{R}\) is Riemann-integrable and vanishes at each point of a dense set. Does it follow that \(\int_0^1 f=0\)? Proof or counterexample, as appropriate.
\item
Let \(f:\mathbb{R}\to\mathbb{R}\) be continuous. Prove that there exists a sequence \((p_n)\) of polynomials that converges to~\(f\) both
\begin{examsubparts}
\item[\textnormal{(i)}]
pointwise on~\(\mathbb{R}\)
\item[\textnormal{(ii)}]
uniformly on each compact set in~\(\mathbb{R}\) (simultaneously).
\end{examsubparts}
\item
Let \(\mathcal{F}\subseteq C(X)\) be equicontinuous and let~\(B\) be the set comprising all points of~\(X\) at which~\(\mathcal{F}\) is bounded; that is, \(a\in B\) exactly when there exists~\(K\) such that \(|f(a)|<K\) whenever \(f\in\mathcal{F}\). Prove that \(B\subseteq X\) is both closed and open.
\item
Let \((f_n)\) be a sequence of measurable real-valued functions on~\(\Omega\). Prove that each of the following subsets of~\(\Omega\) is measurable:
\begin{examsubparts}
\item[\textnormal{(i)}]
\(P=\{\omega:(f_n(\omega))\text{ converges to an irrational number}\}\)
\item[\textnormal{(ii)}]
\(Q=\{\omega:(f_n(\omega))\text{ converges to a rational number}\}\)
\item[\textnormal{(iii)}]
\(R=\{\omega:(f_n(\omega))\text{ converges to a real number}\}\).
\end{examsubparts}
\item
Let \((f_n)\) be a sequence of bounded, measurable, real-valued functions on~\(\Omega\) and assume that \(f_n\to f\) uniformly. Prove that if \(\mu(\Omega)<\infty\) then 
\[
\int_\Omega f_n\,d\mu\to\int_\Omega f\,d\mu
\]
 and show by example that this can fail if \(\mu(\Omega)=\infty\).
\end{examproblems}
\end{document}
